\subsection{Macrospin Effective Fields}

\paragraph{Field convention.}
For a fixed-length macrospin with normalized magnetization $\mvec$, let $w$
denote the magnetic energy density in $\mathrm{J\,m^{-3}}$.  The effective
magnetic field and the corresponding effective flux density are~\cite{ee237-week5}
\begin{equation}
    \Heff
    =-\frac{1}{\mu_0\Ms}\frac{\partial w}{\partial\mvec},
    \qquad
    \vect{B}_{\mathrm{eff}}=\mu_0\Heff.
    \label{eq:macrospin-effective-field}
\end{equation}
Thus, an LLG equation written with $\gamma_0=\mu_0\abs{\gamma}$ uses
$\Heff$ in $\mathrm{A\,m^{-1}}$, whereas one written with
$\abs{\gamma}$ uses $\vect{B}_{\mathrm{eff}}$ in tesla.  The two
conventions are physically equivalent but their field variables must not be
mixed.

Unless stated otherwise, $K_u$ and $A_{12}$ are energy densities in
$\mathrm{J\,m^{-3}}$.  The interfacial $K_i$ in the combination
$K_{\mathrm{bulk}}+K_i/t_{\mathrm{FM}}$ is instead in
$\mathrm{J\,m^{-2}}$, and every saturation magnetization is in
$\mathrm{A\,m^{-1}}$.  The tensors $\mathsf{N}$ and $\mathsf{K}_{fp}$ are
dimensionless.

\subsubsection{One-Sublattice Model}

\paragraph{Energy density and field decomposition.}
For uniaxial anisotropy with easy axis $\unitvect{u}$, a dimensionless
symmetric demagnetization tensor $\mathsf{N}$, an applied field
$\vect{H}_{\mathrm{ext}}$, and a dipolar field $\vect{H}_{\mathrm{dip}}$
produced by another magnetic body, use the uniaxial-anisotropy convention
from slides 13--14 of the EE237 Week 5 lecture notes~\cite{ee237-week5} and
take
\begin{equation}
    w_1
    =K_u\left[1-\bigl(\mvec\cdot\unitvect{u}\bigr)^2\right]
    +\frac{\mu_0\Ms^2}{2}\mvec\cdot\mathsf{N}\mvec
    -\mu_0\Ms\mvec\cdot
    \bigl(\vect{H}_{\mathrm{ext}}+\vect{H}_{\mathrm{dip}}\bigr).
    \label{eq:one-sublattice-energy}
\end{equation}

\paragraph{External field.}
The part of the final term in Eq.~\eqref{eq:one-sublattice-energy} that is
proportional to $\vect{H}_{\mathrm{ext}}$ is the Zeeman coupling to the
applied field.  Its derivative contributes $\vect{H}_{\mathrm{ext}}$ directly
to $\Heff$~\cite{ee237-week5}.

\paragraph{Uniaxial anisotropy field.}
Equation~\eqref{eq:macrospin-effective-field} gives~\cite{ee237-week5}
\begin{equation}
    \vect{H}_{\mathrm{an}}
    =\frac{2K_u}{\mu_0\Ms}
    \bigl(\mvec\cdot\unitvect{u}\bigr)\unitvect{u},
    \qquad
    \vect{B}_{\mathrm{an}}
    =\frac{2K_u}{\Ms}
    \bigl(\mvec\cdot\unitvect{u}\bigr)\unitvect{u}.
    \label{eq:macrospin-anisotropy-field}
\end{equation}
For $\unitvect{u}=\unitvect{z}$, the anisotropy term reduces to
$K_u(1-m_z^2)=K_u\sin^2\theta$, as on slide 13. Near $m_z=1$, slide 14
therefore gives $H_{\mathrm{an},z}=2K_u/(\mu_0\Ms)$.
For a thin film with bulk and interfacial uniaxial anisotropy,
$K_u=K_{\mathrm{bulk}}+K_i/t_{\mathrm{FM}}$.

\paragraph{Demagnetization field.}
Slide 21 of the same notes writes the demagnetization self-energy
as~\cite{ee237-week5}
\begin{align}
    E_{\mathrm{demag}}
     & =-\frac{\mu_0}{2}\int_V
    \vect{M}(\vect{r})\cdot\vect{H}_{\mathrm{demag}}(\vect{r})
    \,\dd^3r,
    \label{eq:macrospin-demag-energy-integral}                 \\
    w_{\mathrm{demag}}
     & =-\frac{\mu_0}{2}\vect{M}\cdot\vect{H}_{\mathrm{demag}}
    =\frac{\mu_0\Ms^2}{2}\mvec\cdot\mathsf{N}\mvec
    \quad\text{(uniform macrospin)}.
    \label{eq:macrospin-demag-energy-density}
\end{align}
The factor $1/2$ avoids double counting the magnet's self-interaction.
Equation~\eqref{eq:macrospin-effective-field} gives
\begin{equation}
    \vect{H}_{\mathrm{demag}}
    =-\Ms\mathsf{N}\mvec,
    \qquad
    \vect{B}_{\mathrm{demag}}
    =-\mu_0\Ms\mathsf{N}\mvec.
    \label{eq:macrospin-demag-field}
\end{equation}
For principal axes,
$\mathsf{N}=\operatorname{diag}(N_x,N_y,N_z)$ and
$N_x+N_y+N_z=1$.

\paragraph{Dipolar field.}
Let a pinned macrospin have volume $V_p$, saturation magnetization
$M_{\mathrm{s},p}$, and direction $\mvec_p$, so its magnetic moment is
$\vect{\mu}_p=V_pM_{\mathrm{s},p}\mvec_p$.
In the point-dipole approximation, at displacement
$\vect{r}=r\unitvect{r}$ from the pinned layer to the free layer, slide 22
gives the dipole formula~\cite{ee237-week5}.  The slide denotes it by
$\vect{H}_i$, but its factor $\mu_0$ makes it a flux density in tesla;
therefore, in the convention of this document,
\begin{align}
    \vect{B}_{\mathrm{dip}}
     & =\frac{\mu_0}{4\pi r^3}
    \left[3\bigl(\vect{\mu}_p\cdot\unitvect{r}\bigr)\unitvect{r}
        -\vect{\mu}_p\right]
    =\mu_0\vect{H}_{\mathrm{dip}},
    \label{eq:macrospin-dipole-flux-density}
\end{align}

\paragraph{Thermal field.}
Following Eq.~(4) of Zhu et al.~\cite{zhu2019voltage}, write the sampled
thermal field as
\begin{align}
    \vect{B}_{\mathrm{th}}
    =\mu_0\Hth
     & =\mathcal{N}_1(0,u_B)\unitvect{x}
    +\mathcal{N}_2(0,u_B)\unitvect{y}
    +\mathcal{N}_3(0,u_B)\unitvect{z},
    \label{eq:thermal-field-discrete}              \\
    u_B
     & =\sqrt{\frac{2\alpha\kB T}
            {V_{\mathrm{FL}}\Ms\abs{\gamma}
                \left(1+\alpha^2\right)\Delta t}}.
    \label{eq:thermal-field-standard-deviation}
\end{align}
Here the $\mathcal{N}_i(0,u_B)$ are independent normal random variables with
zero mean and standard deviation $u_B$, $V_{\mathrm{FL}}$ is the free-layer
volume, and $\Delta t$ is the duration over which a sampled field is held
constant.  With $\abs{\gamma}$ in $\mathrm{rad\,s^{-1}\,T^{-1}}$ and $\Ms$ in
$\mathrm{A\,m^{-1}}$, $u_B$ is in tesla; the corresponding $\vect{H}$-field
standard deviation is $u_H=u_B/\mu_0$.

The source uses independent samples over
$\Delta t=5\,\mathrm{ps}$~\cite{zhu2019voltage}.  When adopting another
integration interval or LLG convention, recompute the noise amplitude
consistently rather than reusing that numerical value.

\paragraph{Total field.}
The total one-sublattice field, with the stochastic field added after taking
the energy derivative, is~\cite{ode-llg-code}
\begin{equation}
    \Heff
    =\vect{H}_{\mathrm{ext}}
    +\vect{H}_{\mathrm{an}}
    +\vect{H}_{\mathrm{demag}}
    +\vect{H}_{\mathrm{dip}}
    +\Hth.
    \label{eq:one-sublattice-total-field}
\end{equation}
Spin-transfer and spin--orbit torques are nonconservative and are not part of
the energy derivative in Eq.~\eqref{eq:macrospin-effective-field}.

\paragraph{Implementation mapping.}
The current one-sublattice \texttt{ODE\_LLG\_integral} implementation stores
flux-density quantities in tesla~\cite{ode-llg-code}.  Its
\texttt{field\_eta.m} evaluates the following expressions, where its input
named \texttt{Hk} is $B_k$:
\begin{align*}
    \vect{B}_{\mathrm{an}}
     & =B_k m_y\unitvect{y} \quad\text{(IMA)},
     &
    \vect{B}_{\mathrm{an}}
     & =B_k m_z\unitvect{z} \quad\text{(PMA)},      \\
    \vect{B}_{\mathrm{demag}}
     & =-\mu_0\Ms\mathsf{N}\mvec,
     &
    \vect{B}_{\mathrm{dip}}
     & =\mu_0\mathsf{K}_{12}\Ms\mvec_{\mathrm{PL}},
\end{align*}
after converting the input $\Ms$ from $\mathrm{emu\,cm^{-3}}$ to
$\mathrm{A\,m^{-1}}$.  It sums these with the applied and thermal fields.

\subsubsection{Two-Sublattice Collinear Model}

\paragraph{Inter-sublattice exchange field.}
Let $\mvec_1$ and $\mvec_2$ be separately normalized sublattice
magnetizations with saturation magnetizations $M_{\mathrm{s},1}$ and
$M_{\mathrm{s},2}$.  The exchange-energy density convention used by the
two-sublattice ferrimagnetic macrospin model of Yuan et
al.~\cite{yuan2023thermal} is
\begin{equation}
    w_{\mathrm{ex},12}=A_{12}\mvec_1\cdot\mvec_2,
    \qquad A_{12}>0,
    \label{eq:two-sublattice-exchange-energy}
\end{equation}
Applying
$\vect{H}_{\mathrm{ex},i}=-(\mu_0M_{\mathrm{s},i})^{-1}
    \partial w_{\mathrm{ex},12}/\partial\mvec_i$ yields, for $j\ne i$,
\begin{equation}
    \vect{H}_{\mathrm{ex},i}
    =-\frac{A_{12}}{\mu_0M_{\mathrm{s},i}}\mvec_j,
    \qquad
    \vect{B}_{\mathrm{ex},i}
    =-\frac{A_{12}}{M_{\mathrm{s},i}}\mvec_j.
    \label{eq:two-sublattice-exchange-field}
\end{equation}

Only the exchange contribution unique to the two-sublattice model is stated
here; the shared external, anisotropy, demagnetization, dipolar, and thermal
field formulas are defined above and are not repeated.  This two-sublattice
exchange extension is publication-grounded and is not implemented by the
\texttt{ODE\_LLG\_integral} repository.
